% =====================================================================
%  CLASE 09 -- Drenaje longitudinal de caminos y carreteras
% =====================================================================
\clase{09}{Drenaje longitudinal de caminos y carreteras}

\section{Drenaje superficial de caminos y carreteras}
Se debe diseñar un correcto drenaje superficial del agua que fluye por los
caminos y carreteras porque:
\begin{itemize}
  \item El agua afecta al pavimento provocando problemas de socavación y
        otros de mecánica de suelos.
  \item No se deben interrumpir cursos de agua existentes, como quebradas.
\end{itemize}
\paragraph{Lectura de referencia:} Manual de Carreteras (MOP), Volumen N°~3
\emph{Instrucciones y Criterios de Diseño} (Capítulo 3.700: Diseño del
drenaje, saneamiento, mecánica e hidráulica fluvial) y Volumen N°~2
\emph{Procedimientos de Estudios Viales} (Capítulo 2.400), disponibles en
\url{https://mc.mop.gob.cl/}.

Existen dos tipos de drenaje:
\begin{itemize}
  \item Drenaje longitudinal.
  \item Drenaje transversal.
\end{itemize}

\section{Drenaje longitudinal}
Las obras hidráulicas son el \textbf{contrafoso} (a media ladera, intercepta
el agua que escurre por la ladera) y la \textbf{cuneta} (junto a la calzada).

\begin{figure}[H]
\centering
\begin{tikzpicture}[font=\small]
  % ladera en planta
  \fill[suelo!80] (0,0) rectangle (7,4);
  \fill[aguaclara] (4.2,0) rectangle (4.6,4);
  \fill[aguaclara] (6.9,0) rectangle (7.05,4);
  \fill[gray!50] (7.05,0) rectangle (8.2,4);
  \foreach \y in {0.4,1.2,2.0,2.8,3.6} \draw[white,very thick] (7.62,\y) -- (7.62,\y+0.4);
  \foreach \x/\l in {0.8/h_7,1.7/h_6,2.6/h_5,3.5/h_4,5.3/h_3,6.2/h_2}
    \draw (\x,0) .. controls (\x+0.2,1.3) and (\x-0.2,2.7) .. (\x+0.1,4) node[pos=0,below]{$\l$};
  \node[below] at (6.9,0) {$h_1$};
  \draw (0.5,2.0) ellipse (0.3 and 0.8); \node at (0.5,2.0) {$h_9$};
  \foreach \x in {1.1,2.0,2.9,3.7,4.9,5.8} \draw[-{Stealth},cyan!60!blue,line width=2pt] (\x,2) -- (\x+0.5,2);
  \draw[flecha,azulusm] (4.4,3) -- (4.4,1);
  \draw[flecha,azulusm] (6.97,3) -- (6.97,1);
  \node[above,azulusm] at (1.5,4.1) {Curvas de nivel};
  \node[above,azulusm] at (4.4,4.1) {Contrafoso};
  \node[above,azulusm] at (7.1,4.1) {Cuneta};
  \node[azulusm,right] at (8.3,0.2) {Vista en planta};
\end{tikzpicture}
\caption{Contrafoso y cuneta en vista en planta: el agua de la ladera es
interceptada antes de llegar a la calzada.}
\end{figure}

\paragraph{Objetivo:} evitar que las aguas que escurren por una ladera
lleguen a la plataforma vial (carretera). El agua que escurre por el
contrafoso y la cuneta se descarga a una quebrada, río, lago o mar.

\begin{figure}[H]
\centering
\begin{tikzpicture}[font=\small]
  \fill[left color=suelo!90!black,right color=suelo!60] (0,0) -- (0,4) -- (1,3.8) -- (1.8,3.5) -- (2.4,3.3)
     -- (2.4,2.6) -- (3.0,2.6) -- (3.1,3.2) -- (3.6,3.0) -- (4.2,2.4) -- (4.8,1.6) -- (5.3,0.8) -- (5.6,0.3) -- (5.6,0) -- cycle;
  \fill[gray!50] (5.6,0) -- (5.6,0.4) -- (5.75,0.4) -- (6.0,0.05) -- (7.6,0.35) -- (8.6,0.25) -- (8.6,0) -- cycle;
  \draw[dashed] (6.0,0.05) -- (8.4,0.05);
  \node at (7.3,0.15) {\scriptsize $i_2$};
  \draw[<-] (2.7,2.6) -- (3.3,3.8) node[above,azulusm]{Contrafoso};
  \draw[<-] (5.85,0.15) -- (5.6,1.3) node[above,azulusm]{Cuneta};
  \draw[<-] (7.6,0.4) -- (8.0,1.1) node[above,azulusm]{Bombeo};
  \node[azulusm] at (6.5,4.5) {Perfil transversal};
\end{tikzpicture}
\caption{Perfil transversal de un camino a media ladera.}
\end{figure}

\subsection{Recomendaciones de diseño de la cuneta}
\begin{itemize}
  \item Ancho: $w\geq0.5$ [m].
  \item Pendiente transversal:
  \begin{itemize}
    \item $i_1\geq8\%$ para cunetas con ancho superior a 0.5 [m].
    \item 30\% máximo para cunetas de ancho de 0.5 [m], para aprovechar de
          mejor forma la capacidad de la cuneta y la eficiencia de los sumideros.
  \end{itemize}
  \item Pendiente longitudinal mínima de la cuneta:
  \begin{itemize}
    \item 0.12\% para cunetas revestidas.
    \item 0.25\% para cunetas sin revestir.
  \end{itemize}
  \item Bombeo o pendiente de la calzada: $i_2\approx2\sim4\%$.
\end{itemize}

\begin{figure}[H]
\centering
\begin{tikzpicture}[font=\small]
  \fill[gray!50] (0,0) rectangle (0.5,2);
  \fill[gray!40] (0.5,0) -- (0.5,0.1) -- (2.2,1.0) -- (6.5,1.5) -- (7.5,1.35) -- (7.5,0) -- cycle;
  \draw (0.5,0.1) -- (2.2,1.0) -- (6.5,1.5) -- (7.5,1.35);
  \draw[dashed] (2.2,0) -- (2.2,2.2);
  \draw[cota] (0.5,1.8) -- (2.2,1.8) node[midway,above]{$w$};
  \draw[dashed] (2.2,1.0) -- (6.8,1.0);
  \draw (1.3,0.1) -- (1.8,0.1); \draw (1.3,0.1) arc (0:28:0.8); \node at (1.5,0.4) {$i_1$};
  \draw (5.8,1.0) arc (180:174:1.2); \node at (5.4,1.15) {$i_2$};
  \draw[<-] (6.5,1.55) -- (6.9,2.1) node[above,azulusm]{Bombeo};
\end{tikzpicture}
\caption{Geometría de la cuneta.}
\end{figure}

Para el correcto diseño del contrafoso y cuneta se debe considerar el cálculo de:
\begin{itemize}
  \item Caudal de diseño.
  \item Altura de agua (velocidad de escurrimiento).
  \item Periodo de retorno.
\end{itemize}

\section{Caudal máximo}
Cuando no se cuenta con registros de caudales, se deben usar fórmulas de
precipitación--escorrentía. En general se trabaja con cuencas pequeñas
(área $<10$ [km$^2$]), lo que se traduce en tiempos de concentración cortos:
\begin{center}
\colorbox{azulusm}{\parbox{0.5\textwidth}{\centering\color{white}\large Duración tormenta $(t_d)>t_c$}}
\end{center}
Por lo tanto, es válido el uso de la \textbf{fórmula racional} (área $<25$ [km$^2$]):
\begin{formula}
$\displaystyle Q_{max}=\frac{C\cdot\bar\imath(t_c)\cdot A}{3.6}\ \ [\text{m}^3/\text{s}]$\\[4pt]
$\bar\imath(t_c)$ en mm/h;\quad $A$: área aportante pluvial en km$^2$
\end{formula}
Para el cálculo del tiempo de concentración se pueden utilizar las fórmulas
vistas en hidrología, pero se debe verificar $t_c\geq10$ [min].

\section{Tiempo de concentración}
Recomendaciones del Volumen 3, Manual de Carreteras:
\begin{table}[H]
\centering\small
\resizebox{\linewidth}{!}{%
\begin{tabular}{lll}
\toprule
\textbf{Autor} & \textbf{Expresión} & \textbf{Observaciones}\\
\midrule
Normas Españolas & $T_c=18\,L^{0.76}/S^{0.19}$ & \\
California Culverts Practice (1942) & $T_c=57\,(L^3/H)^{0.385}$ & Cuencas de montaña\\
Giandotti & $T_c=60\,\dfrac{4A^{0.5}+1.5L}{0.8\,H_m^{0.5}}$ & Cuencas pequeñas con pendiente\\
SCS (1975) & $T_c=3.42\,L^{0.8}\left(\frac{1000}{CN}-9\right)^{0.7}/S^{0.5}$ & Cuencas rurales\\
\bottomrule
\end{tabular}}
\caption{Tabla 3.702.501.A: tiempos de concentración para cuencas.}
\end{table}
Notación: $T_c$: tiempo de concentración (min); $L$: longitud del cauce
(km); $S$: pendiente (\%); $A$: área de la cuenca (km$^2$); $H_m$: diferencia
de nivel (m) entre la cota media de la cuenca y la salida; $H$: diferencia de
nivel total entre cotas extremas de la cuenca (m); $CN$: curva número.

\begin{table}[H]
\centering\small
\resizebox{\linewidth}{!}{%
\begin{tabular}{lll}
\toprule
\textbf{Autor} & \textbf{Expresión} & \textbf{Observaciones}\\
\midrule
Federal Aviation Agency (1970) & $T_c=3.26\,(1.1-C)\,L^{0.5}/(100S)^{0.33}$ & Aeropuertos\\
Izzard (1946) & $T_c=525.28\,(0.0000276\,i+C)\,L_s^{0.33}/(i^{0.667}S^{0.333})$ & Experimentos de laboratorio\\
Morgali y Linsley (1965) & $T_c=7\,L_s^{0.6}\,n^{0.6}/(i^{0.4}S^{0.3})$ & Flujo superficial\\
\bottomrule
\end{tabular}}
\caption{Tabla 3.702.501.B: tiempos de concentración para áreas planas.}
\end{table}
Notación: $L_s$: longitud de escurrimiento superficial (m); $i$: intensidad
de lluvia (mm/h); $C$: coeficiente de escurrimiento; $n$: rugosidad
superficial de Manning. ¡Revisar clases de hidrología!

\section{Precipitación de diseño}
\begin{itemize}
  \item Si se tienen datos horarios se pueden obtener las curvas IDF.
  \item Si sólo se cuenta con registros diarios, se pueden utilizar:
  \begin{itemize}
    \item La fórmula de Grunsky:
    \begin{formula}
    $\displaystyle i_t=i_{24}\sqrt{\frac{24}{t}}\ \ \left[\frac{\text{mm}}{\text{h}}\right]$
    \end{formula}
    \item Coeficientes de duración y frecuencia (V3, Manual de Carreteras):
    \begin{formula}
    $P_t^T=K\cdot CD_t\cdot CF_T\cdot P_D^{10}$
    \end{formula}
  \end{itemize}
\end{itemize}
donde $P_t^T$: lluvia con período de retorno de $T$ años y duración $t$
horas; $P_D^{10}$: lluvia diaria (8 AM a 8 AM) con 10 años de período de
retorno obtenida de una estación pluviométrica; $CD_t$: coeficiente de
duración para $t$ horas; $CF_T$: coeficiente de frecuencia para $T$ años de
período de retorno; $K$: coeficiente de corrección para la lluvia máxima
$P_D^{10}$ medida entre 8 AM y 8 AM respecto de las 24 h más lluviosas de la
tormenta, para el que se ha adoptado un valor $K=1.1$.

Para duraciones inferiores a una hora:
\begin{itemize}
  \item La fórmula de Grunsky.
  \item La fórmula de Bell (1969, V3 Manual de Carreteras):
  \begin{formula}
  $P_t^T=\left(0.54\,t^{0.25}-0.50\right)\left(0.21\ln T+0.52\right)P_1^{10}$
  \end{formula}
  con $P_t^T$: precipitación (mm) con período de retorno $T$ años y duración
  $t$ minutos; $P_1^{10}$: precipitación (mm) con $T=10$ años y duración de 1 hora.
\end{itemize}

\section{Coeficiente de escorrentía}
\begin{table}[H]
\centering
\begin{tabular}{lc}
\toprule
\textbf{Tipo de terreno} & \textbf{Coeficiente de escurrimiento}\\
\midrule
Pavimentos de adoquín & 0.50 -- 0.70\\
Pavimentos asfálticos & 0.70 -- 0.95\\
Pavimentos en concreto & 0.80 -- 0.95\\
Suelo arenoso con vegetación y pendiente 2\% -- 7\% & 0.15 -- 0.20\\
Suelo arcilloso con pasto y pendiente 2\% -- 7\% & 0.25 -- 0.65\\
Zonas de cultivo & 0.20 -- 0.40\\
\bottomrule
\end{tabular}
\caption{Tabla 3.702.503.A: coeficientes de escurrimiento $C$.}
\end{table}

\begin{table}[H]
\centering\scriptsize
\begin{tabularx}{\textwidth}{>{\raggedright}p{1.8cm}XXXX}
\toprule
\textbf{Factor} & \textbf{Extremo} & \textbf{Alto} & \textbf{Normal} & \textbf{Bajo}\\
\midrule
Relieve & 0.28--0.35: escarpado con pendientes mayores que 30\% & 0.20--0.28: montañoso con pendientes entre 10 y 30\% & 0.14--0.20: con cerros y pendientes entre 5 y 10\% & 0.08--0.14: relativamente plano con pendientes menores al 5\%\\
\addlinespace
Infiltración & 0.12--0.16: suelo rocoso, o arcilloso con capacidad de infiltración despreciable & 0.08--0.12: suelos arcillosos o limosos con baja capacidad de infiltración, mal drenados & 0.06--0.08: normales, bien drenados, textura mediana, limos arenosos, suelos arenosos & 0.04--0.06: suelos profundos de arena u otros suelos bien drenados con alta capacidad de infiltración\\
\addlinespace
Cobertura vegetal & 0.12--0.16: cobertura escasa, terreno sin vegetación o escasa cobertura & 0.08--0.12: poca vegetación, terrenos cultivados o naturales, menos del 20\% del área con buena cobertura vegetal & 0.06--0.08: regular a buena; 50\% del área con praderas o bosques, no más del 50\% cultivado & 0.04--0.06: buena a excelente; 90\% del área con praderas, bosques o cobertura equivalente\\
\addlinespace
Almacenamiento superficial & 0.10--0.12: despreciable, pocas depresiones superficiales, sin zonas húmedas & 0.08--0.10: baja, sistema de cauces superficiales pequeños bien definidos, sin zonas húmedas & 0.06--0.08: normal; posibilidad de almacenamiento buena, zonas húmedas, pantanos, lagunas y lagos & 0.04--0.06: capacidad alta, sistema hidrográfico poco definido, buenas planicies de inundación o gran cantidad de zonas húmedas, lagunas o pantanos\\
\bottomrule
\end{tabularx}
\caption{Tabla 3.702.503.B: coeficientes de escorrentía $C$ para $T=10$ años
(se suman los aportes de cada factor). Si $T>10$ años, amplificar el
resultado por: $T=25$: $C\times1.10$; $T=50$: $C\times1.20$; $T=100$: $C\times1.25$.}
\end{table}

\section{Altura de agua $h$}
Por tratarse de flujo a superficie libre, con el ancho o algún otro dato
geométrico dado se utiliza la fórmula de Manning:
\begin{formula}
$\displaystyle\frac{Q\cdot n}{\sqrt s}=\frac{\Omega^{5/3}}{\Psi^{2/3}}$\qquad
$\Omega$: área de escurrimiento;\quad $\Psi$: perímetro mojado
\end{formula}
\begin{itemize}
  \item Generalmente los contrafosos están construidos de tierra: $n=0.023\sim0.025$.
  \item Se espera que el escurrimiento sea subcrítico; sin embargo, se debe
        verificar con la altura crítica.
\end{itemize}

\subsection{Rugosidad de Manning}
Se pueden usar rugosidades características, dependiendo del tipo de
revestimiento del canal (extracto de la Tabla 3.705.1.A, valores mínimo,
normal y máximo):
\begin{table}[H]
\centering\small
\begin{tabular}{llccc}
\toprule
\textbf{Revestimiento} & \textbf{Descripción} & \textbf{Mín.} & \textbf{Normal} & \textbf{Máx.}\\
\midrule
Hormigón & Platachado & 0.011 & 0.013 & 0.015\\
 & Alisado con regla & 0.013 & 0.015 & 0.016\\
 & Alisado con ripio a la vista en el fondo & 0.015 & 0.017 & 0.020\\
 & Sin alisar & 0.014 & 0.017 & 0.020\\
 & Gunita, sección regular & 0.016 & 0.019 & 0.023\\
 & Gunita, sección ondulada & 0.018 & 0.022 & 0.025\\
Fondo de hormigón alisado & Piedra acomodada en mortero & 0.015 & 0.017 & 0.020\\
\quad con lados de: & Piedra distribuida al azar en mortero & 0.017 & 0.020 & 0.024\\
 & Empedrado o enrocado (rip rap) & 0.020 & 0.030 & 0.035\\
Fondo de grava con lados de: & Hormigón (con moldaje) & 0.017 & 0.020 & 0.025\\
 & Empedrado o rip rap & 0.023 & 0.033 & 0.036\\
Ladrillo & Barnizado o vidriado & 0.011 & 0.013 & 0.015\\
 & En mortero de cemento & 0.012 & 0.015 & 0.018\\
Albañilería & Empedrado cementado & 0.017 & 0.025 & 0.030\\
 & Empedrado libre & 0.023 & 0.032 & 0.035\\
Asfalto & Liso & 0.013 & 0.013 & --\\
 & Rugoso & 0.016 & 0.016 & --\\
Cubierto con vegetación & & 0.030 & -- & 0.500\\
\bottomrule
\end{tabular}
\caption{Coeficientes de rugosidad de Manning en canales.}
\end{table}

\subsection{Cálculo de altura normal y crítica}
Utilizar lo aprendido en Hidráulica Teórica. Se puede ayudar de software
(p.~ej.\ HCANALES) para determinar rápidamente $h_n$ y $h_c$.

\section{Velocidades de escurrimiento}
\paragraph{Velocidades máximas} (evitar problemas de socavación):
\begin{table}[H]
\centering
\begin{tabular}{cc}
\toprule
$V_{max}$ [m/s] & Tipo de suelo\\
\midrule
0.6 & Limo\\
0.9 & Arena\\
1.2 & Arcilla\\
3.0 & Asfalto\\
4.5 & Hormigón\\
\bottomrule
\end{tabular}
\end{table}
\paragraph{Velocidades mínimas} (evitar problemas de sedimentación):
\begin{formula}
$v_{min}=0.4\sim0.5$ [m/s]
\end{formula}
La condición de velocidad mínima también se puede controlar por la pendiente
de diseño (fondo del contrafoso):
\begin{table}[H]
\centering
\begin{tabular}{cc}
\toprule
$S_{min}$ [\%] & Tipo de suelo\\
\midrule
0.25 & Tierra\\
0.15 & Asfalto\\
\bottomrule
\end{tabular}
\end{table}

\paragraph{Revancha de seguridad ($R$):} evitar el rebalse del contrafoso
debido a alguna condición no prevista en el diseño.
\begin{formula}
$R=0.15\sim0.2\,h$
\end{formula}

\section{Geometría y periodo de retorno}
\paragraph{Geometría}
\begin{itemize}
  \item En caso de contrafosos de suelo natural, se recomiendan secciones de
        diseño trapezoidales.
  \item Además, se debe verificar que $\theta<\phi$, con $\phi$: ángulo de
        fricción del suelo.
\end{itemize}

\begin{figure}[H]
\centering
\begin{tikzpicture}[font=\small]
  \fill[suelo!80] (-0.8,0) -- (0.4,0) -- (-0.3,2.6) -- (-0.8,2.6) -- cycle;
  \fill[suelo!80] (4.4,0) -- (5.6,0) -- (5.6,2.6) -- (5.1,2.6) -- cycle;
  \fill[gray!30] (-0.3,2.6) -- (5.1,2.6) -- (4.85,1.7) -- (-0.05,1.7) -- cycle;
  \fill[top color=aguaclara,bottom color=agua] (-0.05,1.7) -- (4.85,1.7) -- (4.4,0) -- (0.4,0) -- cycle;
  \draw[thick] (-0.8,2.6) -- (-0.3,2.6) -- (0.4,0) -- (4.4,0) -- (5.1,2.6) -- (5.6,2.6);
  \pic at (0.4,1.72) {nivelagua};
  \draw[dashed] (4.4,0) -- (5.6,0);
  \draw (5.1,0) arc (0:75:0.7); \node at (5.25,0.35) {$\theta$};
  \draw[cota] (5.9,0) -- (5.9,1.7) node[midway,right]{$h$};
  \draw[cota] (5.9,1.7) -- (5.9,2.6) node[midway,right]{$R=0.15\sim0.2\,h$};
\end{tikzpicture}
\caption{Sección trapezoidal del contrafoso con revancha $R$.}
\end{figure}

\paragraph{Periodo de retorno de diseño}
\begin{itemize}
  \item Dependerá del tipo de camino y su importancia.
  \item También se puede diseñar considerando el riesgo de falla:
\end{itemize}
\begin{formula}
$\displaystyle R=1-\left[1-\frac1T\right]^n$
\end{formula}
\begin{table}[H]
\centering
\begin{tabular}{lccccc}
\toprule
& \multicolumn{2}{c}{$T$ [años]} & Vida útil & \multicolumn{2}{c}{Riesgo de falla $R$ [\%]}\\
\cmidrule(lr){2-3}\cmidrule(lr){5-6}
\textbf{Tipo de obra} & Diseño & Verificación & supuesta $n$ [años] & Diseño & Verificación\\
\midrule
Carreteras & 10 & 25 & 10 & 65 & 34\\
Caminos & 5 & 10 & 5 & 67 & 41\\
\bottomrule
\end{tabular}
\end{table}

\section{Bajadas de agua: rápidos o canal en gradería}
\begin{itemize}
  \item Su función es proteger de la erosión a los taludes de terraplenes y
        cortes, conduciendo el agua hasta un sitio seguro (cuneta al pie del talud).
  \item Diámetro mínimo de 200 mm para $H/V\geq4$.
  \item Gradería para razones $H/V$ 4:1 o menores.
  \item Pueden ser canaletas semicirculares en metal corrugado, $H/V$ 2:1.
  \item Velocidades de 5 m/s para el diseño hidráulico.
\end{itemize}

\begin{figure}[H]
\centering
\begin{tikzpicture}[font=\small]
  \fill[gray!25] (0,0) -- (0,0.4) -- (6,3.4) -- (6,0) -- cycle;
  \draw (0,0.4) -- (6,3.4);
  \draw[thick,azulusm] (0,0.4) \foreach \i in {0,...,6}{ -- ++(0,0.43) -- ++(0.86,0) };
  \foreach \i in {0,...,6} \draw[flecha,cyan!70!blue] ({0.86*\i+0.7},{0.43*\i+0.95}) -- ++(-0.5,-0.25);
  \node[align=center] at (4.5,1.0) {Canal en gradería\\(corte longitudinal)};
\end{tikzpicture}
\caption{Bajada de agua escalonada (canal en gradería) sobre un talud.}
\end{figure}
